\documentclass[11pt]{article}

\usepackage{amsmath, amssymb, mathtools}
\usepackage[margin=1in]{geometry}

\title{Exact and Fast Metric Path-Following Coordinates}
\author{}
\date{}

\begin{document}
\maketitle

\section{Setup}

Let \(Q\) be the configuration manifold of the robot and let
\[
    g_q(v,w) = v^\top M(q) w
\]
be the kinetic-energy Riemannian metric. For the two-link example,
\[
    Q = S^1 \times S^1,
\]
with coordinates \(q = (\theta_1,\theta_2)\), and \(M(q)\) is the manipulator
mass matrix.

Let
\[
    \gamma : [0,1] \to Q
\]
be a non-self-intersecting reference path. The path-following coordinates are
intended to represent a nearby configuration \(q\) by
\[
    (\eta,\xi),
\]
where \(\eta\) is the closest point on the path and \(\xi\) is the transverse
displacement from \(\gamma(\eta)\) to \(q\).

\section{Exact Theoretical Construction}

The exact Riemannian construction uses the geodesic distance induced by \(g\).
The coordinate \(\eta\) is defined by
\[
    \eta
    =
    \operatorname*{argmin}_{s \in [0,1]}
    d_g(\gamma(s),q)^2.
\]
Equivalently, if \(\log^g_p(q)\) denotes the Riemannian logarithm, then
\[
    \eta
    =
    \operatorname*{argmin}_{s \in [0,1]}
    \left\|
        \log^g_{\gamma(s)}(q)
    \right\|^2_{g_{\gamma(s)}}.
\]
The transverse vector is then
\[
    \xi
    =
    \log^g_{\gamma(\eta)}(q).
\]

This is the clean geometric object. It has the property that
\[
    q
    =
    \operatorname{Exp}^g_{\gamma(\eta)}(\xi),
\]
where \(\operatorname{Exp}^g\) is the Riemannian exponential map.

\subsection{Computational Cost}

For a non-constant metric \(M(q)\), the logarithm
\[
    \log^g_p(q)
\]
is not a simple coordinate difference. It is obtained by solving the geodesic
boundary-value problem
\[
    \nabla_{\dot c}\dot c = 0,
    \qquad
    c(0)=p,
    \qquad
    c(1)=q.
\]
In local coordinates this becomes
\[
    \ddot c^k
    +
    \Gamma^k_{ij}(c)\dot c^i \dot c^j
    =
    0,
\]
where the Christoffel symbols \(\Gamma^k_{ij}\) are determined by the metric.

The expensive part is not merely solving one geodesic problem. The closest
point search evaluates the objective for many candidate values of \(s\). Thus
one controller call may require many BVP solves:
\[
    s
    \longmapsto
    \left\|
        \log^g_{\gamma(s)}(q)
    \right\|^2_{g_{\gamma(s)}}.
\]
This is why the exact implementation is too slow for real-time control.

\section{Fast Local Method on a Configuration Submanifold}

Now let \(Q\) be an arbitrary \(d\)-dimensional configuration submanifold,
possibly embedded in an ambient vector space \(\mathcal A \cong \mathbb R^m\).
The fast method does not require \(Q\) to be \(S^1\times S^1\).  What it needs
is a cheap local replacement for the Riemannian logarithm.

For each \(p\in Q\), choose a local retraction
\[
    R_p : U_p \subset T_pQ \to Q,
    \qquad
    R_p(0)=p,
    \qquad
    D R_p(0)=I.
\]
The associated local inverse, when it exists, is denoted
\[
    \ell_p(q) \approx R_p^{-1}(q) \in T_pQ.
\]
This \(\ell_p(q)\) is the fast logarithm.  It is not the Riemannian logarithm
\(\log^g_p(q)\); it is a chart or retraction displacement that is cheap to
evaluate.

Examples include:
\[
    \ell_p(q)
    =
    \chi_p(q)-\chi_p(p)
\]
in a chosen local coordinate chart \(\chi_p\), or, for an embedded
submanifold, an ambient displacement projected onto \(T_pQ\),
\[
    \ell_p(q)
    =
    P_{T_pQ}(q-p),
\]
possibly followed by a projection/retraction back to \(Q\).  For
\(S^1\times S^1\), this is exactly the wrapped joint-space difference.

The fast closest-point coordinate is then
\[
    \eta_{\mathrm{fast}}
    =
    \operatorname*{argmin}_{s\in[0,1]}
    g_{\gamma(s)}
    \left(
        \ell_{\gamma(s)}(q),
        \ell_{\gamma(s)}(q)
    \right).
\]
The fast displacement is
\[
    \xi_{\mathrm{fast}}
    =
    \ell_{\gamma(\eta_{\mathrm{fast}})}(q)
    \in T_{\gamma(\eta_{\mathrm{fast}})}Q.
\]

This keeps the metric \(g\) in the objective but replaces the expensive
geodesic BVP logarithm by a local algebraic displacement.  In a sufficiently
small neighborhood of \(p\),
\[
    d_g(p,q)^2
    =
    g_p(\ell_p(q),\ell_p(q))
    +
    \text{higher-order terms}.
\]
Thus the fast method is the local quadratic model of the Riemannian distance.

\section{Lie Group Specialization}

If the configuration manifold \(Q\) is a Lie group \(G\), there is a natural
choice of retraction:
\[
    R_p(\xi)
    =
    p\,\exp_G(\xi),
    \qquad
    \ell_p(q)
    =
    \log_G(p^{-1}q),
\]
where \(\exp_G\) and \(\log_G\) are the Lie group exponential and logarithm.
The corresponding closest-point coordinate is
\[
    \eta_{\mathrm{Lie}}
    =
    \operatorname*{argmin}_{s\in[0,1]}
    g_{\gamma(s)}
    \left(
        \log_G(\gamma(s)^{-1}q),
        \log_G(\gamma(s)^{-1}q)
    \right).
\]
Then
\[
    \xi_{\mathrm{Lie}}
    =
    \log_G(\gamma(\eta_{\mathrm{Lie}})^{-1}q).
\]

This is still generally not the same as the Riemannian logarithm
\(\log^g_p(q)\) for the kinetic-energy metric.  It becomes the exact
Riemannian logarithm only when the chosen Riemannian metric has compatible
group invariance.  In the two-link example,
\[
    Q = S^1\times S^1,
\]
and the group is abelian.  Therefore
\[
    \log_G(p^{-1}q)
    =
    \operatorname{wrap}(q-p),
\]
component by component.  This is exactly the same displacement currently used
by the fast method.  The implementation exposes it as a separate option
because the geometric interpretation is different, even though the numerical
result is identical for this particular torus.

\section{Metric Tangent and Transverse Frame}

Let \(d=\dim Q\).  Along the path, define the metric-unit tangent
\[
    T(\eta)
    =
    \frac{\gamma'(\eta)}
         {\|\gamma'(\eta)\|_{g_{\gamma(\eta)}}}.
\]
Choose a metric-orthonormal transverse frame
\[
    N(\eta)
    =
    \begin{bmatrix}
        N_1(\eta) & \cdots & N_{d-1}(\eta)
    \end{bmatrix}
\]
satisfying
\[
    g_{\gamma(\eta)}(T,N_a)=0,
    \qquad
    g_{\gamma(\eta)}(N_a,N_b)=\delta_{ab}.
\]
The transverse coordinate is the metric projection of the fast displacement
onto this frame:
\[
    \zeta
    =
    N(\eta)^\flat \xi_{\mathrm{fast}},
    \qquad
    \zeta_a
    =
    g_{\gamma(\eta)}
    \left(
        \xi_{\mathrm{fast}},
        N_a(\eta)
    \right).
\]
Here \(\zeta\in\mathbb R^{d-1}\).  In matrix coordinates, if \(G(\gamma(\eta))\)
is the metric matrix and \(N\) is the matrix whose columns are \(N_a\), then
\[
    \zeta
    =
    N^\top G(\gamma(\eta))\xi_{\mathrm{fast}}.
\]

For the two-link example \(d=2\), so there is only one normal vector and
\(\zeta\) is the scalar \(\xi_\perp\).  If
\[
    \alpha = M(\gamma(\eta))T,
\]
then the implementation chooses the metric-normal direction by taking a
Euclidean perpendicular to \(\alpha\):
\[
    \tilde N =
    \begin{bmatrix}
        -\alpha_2 \\
        \alpha_1
    \end{bmatrix},
    \qquad
    N =
    \frac{\tilde N}{\|\tilde N\|_{M(\gamma(\eta))}}.
\]

\section{Controller}

The corresponding local tubular chart is
\[
    q = \Phi(\eta,\zeta)
      =
      R_{\gamma(\eta)}
      \left(
          N(\eta)\zeta
      \right),
    \qquad
    \zeta\in\mathbb R^{d-1}.
\]
The exact theoretical version replaces the retraction \(R\) by the Riemannian
exponential \(\operatorname{Exp}^g\).  The fast implementation uses the same
formula but chooses \(R\) to be something cheap, such as a wrapped joint chart
or an ambient projection/retraction.

Let
\[
    y =
    \begin{bmatrix}
        \eta \\
        \zeta
    \end{bmatrix}
    \in \mathbb R^d.
\]
In local coordinates on \(Q\), the acceleration relation is
\[
    \dot q = D\Phi(y)\dot y,
    \qquad
    \ddot q =
    D\Phi(y)\ddot y
    +
    \frac{d}{dt}\bigl(D\Phi(y)\bigr)\dot y.
\]
For an embedded representation \(q\in\mathcal A\), the same equation holds in
the ambient space with \(\dot q,\ddot q\in T_qQ\), and \(\dot y\) can be
obtained from a local inverse or metric-weighted pseudoinverse of \(D\Phi(y)\).
The second term is the path-curvature/centripetal term that is missing from a
pure velocity-field controller.

The path coordinate command is
\[
    \ddot \eta
    =
    \ddot \eta_{\mathrm{ref}}
    +
    k_\eta(\dot\eta_{\mathrm{ref}}-\dot\eta),
\]
where \(\dot\eta_{\mathrm{ref}}\) is chosen so
\(\|\gamma'(\eta)\dot\eta_{\mathrm{ref}}\|_{g_{\gamma(\eta)}}\) equals the
desired metric speed.  The transverse command is
\[
    \ddot \zeta
    =
    -K_p \zeta - K_d \dot\zeta.
\]
The commanded configuration acceleration is then
\[
    \ddot q_{\mathrm{cmd}}
    =
    D\Phi(y)
    \begin{bmatrix}
        \ddot\eta \\
        \ddot\zeta
    \end{bmatrix}
    +
    \frac{d}{dt}\bigl(D\Phi(y)\bigr)\dot y.
\]
Finally, the torque is
\[
    \tau
    =
    M(q)\ddot q_{\mathrm{cmd}}
    +
    C(q,\dot q)\dot q
    +
    G(q).
\]

\section{Why the Fast Method is Faster}

The exact method repeatedly performs:
\[
    \text{candidate }s
    \quad\Rightarrow\quad
    \text{solve geodesic BVP}
    \quad\Rightarrow\quad
    \text{evaluate distance}.
\]
The fast method performs:
\[
    \text{candidate }s
    \quad\Rightarrow\quad
    \text{local retraction inverse } \ell_{\gamma(s)}(q)
    \quad\Rightarrow\quad
    \text{metric quadratic form}.
\]

Thus the inner loop changes from nonlinear boundary-value problem solves to
simple vector and matrix operations. The approximation is appropriate when the
state remains in a sufficiently small tube around a non-self-intersecting path.

\section{Limitations}

The fast method is local. It should not be interpreted as the exact Riemannian
logarithm for large displacements. It requires a computable local chart,
projection, or retraction whose inverse is unique in the tube around the path.
It also assumes that the path-following tube has a unique closest branch.
Self-intersecting paths should be avoided because the coordinate \(\eta\) is
not unique at intersections.

For an arbitrary configuration submanifold, the method is therefore available
when the following data are available:
\[
    \gamma(\eta),\qquad
    \gamma'(\eta),\qquad
    g_q,\qquad
    R_p,\qquad
    \ell_p.
\]
The robot dynamics do not need to have the special structure of the two-link
arm, beyond being fully actuated in the coordinates used for inverse dynamics.
The geometry bottleneck is moved from solving geodesic BVPs to evaluating the
chosen local retraction model.

\end{document}
